\bookchapter{第三十二章：スマホ安全救急マニュアル}

\chapterintro{この章は「緊急時対応マニュアル」です。できれば一生使わないでほしいのですが、もしスマホの紛失·盗難·詐欺に遭ったとき、このマニュアルがあれば被害を最小限に抑えられます。地震や大雨などの災害が来たときも、スマホは警報の受信、居場所の伝達、連絡の維持に欠かせない道具です。今のうちに一度通読して、大事な手順をスクリーンショットで保存するか、印刷して家に置いておくのがおすすめです。何かあったときにはスマホがもう手元にないかもしれず、災害時には電池切れや圏外の可能性もあるからです。}

\booksection{シーン0：地震／極端な天候——スマホが命を守るしくみ}

守るべきは命の安全であって、アカウントやお金ではありません。警報の設定方法は第十三章「天気警報·自然災害警報·地震警報」を参照。ここでは警報が鳴ったあとの行動だけを説明します。

\textbf{地震警報が鳴ったら：まず身を伏せてから避難。動画撮影に夢中にならない}\par

\begin{enumerate}
  \item 近くで身を守る：まずしゃがみ、頭を覆い、しっかりつかまる（テーブルの下、耐力壁の隅、頑丈な家具のそば）。頭と首を保護し、ガラス·外壁·シャンデリアから離れます。
  \item むやみに建物の外へ走り出さない：揺れの本体は多くの場合十数秒から数十秒で、階段のほうがかえって危険なことも。エレベーターは絶対に使わない。
  \item 揺れが収まってから落ち着いて避難：階段を使い、余震と落下物に注意。開けた場所へ移動し、電柱·看板·倒壊しそうな建物から離れる。
  \item スマホは本当に必要な場面で：緊急連絡先へ短いメッセージを送る（「私は○○にいます。無事です／負傷しています。助けが必要です」）ほうが、長電話より電池がもちます。送れる状況なら、まずはテキスト（SMS／WeChat）で。
  \item 圏外のとき：第七章の衛星SOS／衛星メッセージ（対応機種の場合）を参照。または高い場所·開けた場所へ移動してから再試行。状況によっては、ネット不要の Bluetooth トランシーバーも役立ちます。
  \item 不要な機能を切る：緊急時は機内モードを入れて定期的に解除、Bluetooth のスキャンをオフ、画面の明るさを下げる——だけで数時間は長持ちします。最後の電力は、安否連絡と救援情報のために取っておきましょう。
\end{enumerate}

\begin{tcbwarn}{大雨／台風／山あいの洪水：警報の段階で動き出す}
\begin{itemize}
  \item 赤色警報を見たら：外出はできるだけ控え、地下空間·アンダーパス·河川·急な斜面から離れる。
  \item 屋外で取り残されたら：頑丈な建物の上の階へ移動し、無理に水の中へ入らない（冠水した場所にはマンホールの蓋がないところや、漏電の危険がある）。119／110 にかけるときは、目印と水位をはっきり伝える。
  \item 出かける前：天気アプリと地域の警報を流す公式アカウントで危険を確認。長距離ドライブでは、山沿い·水辺·崖沿いの区間を避ける。
  \item 家族での取り決め：「連絡が取れなくなったら、最後の WeChat／SMS の位置情報を正とする」と決めておく。全員が必死に発信し続けて電池を使い切るのを防げます。
\end{itemize}

\end{tcbwarn}

\begin{tcbpractice}{ふだんから済ませておく4つのこと}
\begin{enumerate}
  \item システムの地震警報や正規の警報アプリを有効にし、マナーモードでも鳴るよう許可しておく。
  \item 連絡先の先頭に「緊急連絡先」を2〜3人置き、ロック画面に表示できるようにする（iOS／Android のメディカルIDやロック画面情報）。
  \item 家とカバンにモバイルバッテリーを用意。長距離移動の前には満充電してケーブルも持参。
  \item 連絡帳を見なくてもかけられる番号（家族の携帯番号など）を最低1つ覚え、紙に書いて財布に入れる。
\end{enumerate}

\end{tcbpractice}

\booksection{シーン1：スマホを紛失した／盗まれた——ゴールデン30分}

スマホを紛失してからの最初の30分は「ゴールデンタイム」です。適切に対処すれば、お金とプライバシーを守れます。以下の順に、一歩も飛ばさず実行してください：

\textbf{紛失後30分の標準的な手順}\par

\begin{enumerate}
  \item 人のスマホを借りて自分の番号に発信（最初の1分）：近くに落ちていて、親切な人が取ってくれれば取り返せるかもしれません。
  \item スマホを遠隔ロック（2〜5分目）：人のスマホかパソコンから「探す」系サービスにログインします： iPhone：iCloud.com ＞ 探す ＞ 自分のデバイスを選択 ＞ 紛失としてマーク
  \item Android の各ブランド：HUAWEI は cloud.huawei.com／Xiaomi は i.mi.com／OPPO は cloud.oppo.com など
  \item SIM カードの利用停止（5〜10分目）：キャリアのサポート窓口（移動10086／聯通10010／電信10000）に電話し、スマホを紛失したと伝えて SIM カードのサービス停止を依頼します。ここが最重要——詐欺師があなたの SIM カードを手にすると、認証コードを受け取られ、すべてのパスワードをリセットされてしまいます。
  \item 決済アカウントの凍結（10〜20分目）： Alipay：95188 で利用停止、または人のスマホで Alipay にログイン ＞ マイ ＞ 設定 ＞ セキュリティセンター ＞ 緊急凍結
  \item WeChat：110 に通報したうえで WeChat サポートに連絡、または WeChat セキュリティセンター ＞ アカウントを凍結
  \item 銀行アプリ：キャッシュカード裏面のサポート電話にかけ、ネットバンキングの凍結を依頼
  \item 知人·家族への連絡（20〜25分目）：モーメンツで知らせるか、家族に拡散を頼み、スマホを紛失したことを伝えて、詐欺師があなたになりすまして他人を騙すのを防ぎます。
  \item 警察への通報（25〜30分目）：110 に電話するか、最寄りの派出所で被害を届け出て、受理票を受け取ります。この後の保険請求や銀行とのやり取りでも、この受理票が必要になります。
\end{enumerate}

\booksection{シーン2：詐欺に騙された——72時間の勝負期間}

騙されたと気づいたら、いちばん大事なのは冷静さです。慌てて間違った判断をすると、二次被害につながります。以下の手順で進めてください：

\begin{enumerate}
  \item まずすべての送金を即座に中止：詐欺師が何と言おうと、1円も送ってはいけません。「もう一度保証金を送れば前のお金が戻る」は、典型的な二次詐欺の口上です。
  \item 証拠をすべて残す：チャット履歴を消さない！通話履歴を消さない！チャットのスクリーンショット、送金記録、相手のアカウント情報を保存します。これらは通報と資金取り戻しのカギとなる証拠です。
  \item すぐに96110へ電話：全国の詐欺被害相談専用ダイヤルです。オペレーターが緊急の支払い停止を案内してくれます——送金から間もない場合（通常24時間以内）なら、銀行が相手の口座を凍結して資金を止められる可能性があります。
  \item 110 に通報：スマホ、身分証、送金の証憑を持って最寄りの派出所へ行き、事情聴取を受けます。
  \item 銀行／決済サービスに連絡：被害状況を伝え、関与した口座の凍結を依頼します。クレジットカードを不正利用された場合は、チャージバック（支払い取消し）の手続きを求めましょう。
  \item パスワードをすべて変更：別の安全なデバイスで、Apple ID／各ブランドのアカウント、WeChat、Alipay、銀行アプリのパスワードを変更し、二段階認証を有効にします。
\end{enumerate}

\booksection{シーン3：スマホがウイルス／トロイの木馬に感染したら——どうする？}

次のような症状が出たら——理由もなく広告がポップアップする、知らないアプリが勝手に入る、通話料が急増する、データ通信量が異常に減る——マルウェアに感染した可能性が高いです：

\begin{itemize}
  \item ステップ1：セーフモードに入る——Android は電源ボタンを長押しして電源を切り、もう一度長押しして起動。起動時のバイブレーションが来たら音量下ボタンを押し続け、セーフモードに入るまで待ちます（画面左下に「セーフモード」と表示）。セーフモードではサードパーティ製アプリがすべて無効になるので、原因の切り分けに便利です。
  \item ステップ2：怪しいアプリを削除——セーフモードで［設定］＞［アプリ管理］を開き、インストール日時順に並べて、最近入れた心当たりのないアプリを見つけて削除します。
  \item ステップ3：セキュリティアプリでスキャン——「Tencent Mobile Manager」や「360 Mobile Guard」をダウンロードして、フルスキャンを実行します。
  \item ステップ4：初期化（最終手段）——ここまでで解決しない場合は、大事なデータをバックアップしてから工場出荷状態に戻します。注意：初期化するとすべてのデータが消えるので、写真·連絡先などの重要なファイルは必ず先にバックアップしてください。
  \item iPhone は基本的にウイルスに感染しません——iOS のサンドボックス機構のおかげでマルウェアが生き残りにくいためです。iPhone に異常がある場合、ほとんどはアプリのバグか設定の問題で、ウイルスではありません。初期化すれば99\%の問題は解決します。
\end{itemize}

\booksection{シーン4：アカウントを乗っ取られた——素早く主導権を取り戻す}

\begin{enumerate}
  \item すぐにパスワードを変更：まだログインできるなら、真っ先に変えます。新しいパスワードは旧パスワードとまったく別のものに。
  \item ほかのデバイスを締め出す：アカウントのセキュリティ設定で「ログイン端末の管理」を探し、知らないデバイスをすべて強制ログアウトさせます。
  \item ログイン履歴を確認：直近7日間のログイン記録を見て、知らない場所からのログインがあれば、アカウントは間違いなく盗まれています。
  \item 二段階認証を有効にする：まだオンにしていないなら、今すぐオンに。アカウントの再被害を防ぐいちばん効果的な方法です。
  \item 連携情報を確認：アカウントに紐づく電話番号やメールアドレスが勝手に変わっていないか確認します。変わっていたら、アカウントの復旧申し立て機能で取り戻しましょう。
  \item 決済設定を確認：被害が WeChat／Alipay なら、パスワードなし決済や自動引き落としの設定が変更されていないか、銀行カードが見知らぬアカウントに追加されていないかを確認します。
\end{enumerate}

\booksection{シーン5：脅迫メッセージが届いた——怖がらない、払わない}

よくある脅迫の手口は2つです。(1)「あなたのわいせつ動画／写真を手に入れた。金を払わなければ連絡先全員に送る」。(2)「あなたのスマホはハッカーに支配された。金を払わなければロックする」。こうした状況に遭ったら：

\begin{itemize}
  \item 99.9\%は詐欺です：相手の手元にあなたのプライバシーデータは何もありません。同じ脅迫文を何万人かに一斉送信して、怖がって払う人がいればもうけ、という仕組みです。
  \item 返信しない、お金を払わない：返信は「この番号は生きている」と相手に教えるのと同じで、「脅せそうなターゲット」の印をつけられたも同然。その後の嫌がらせが増えます。
  \item スクリーンショットで証拠を保存：そのうえで相手をブロックします。
  \item パスワードを変更：心配なら、メインのアカウントのパスワードを変えておけば安心です。
  \item 本当にプライバシー漏えいの危険がありそうなら：警察に通報を。警察にはこの種の事件を扱う経験があります。
\end{itemize}

\booksection{シーン6：AI 顔すり替え／AI 声まね詐欺——見たものが真実ではなくなる時代}

2024年以降、AI による顔すり替えと声のまね技術が進み、詐欺師がビデオ通話の中であなたの家族·上司·友人の顔や声に「なりすます」ことができるようになりました。「ママ、トラブルに遭ったから、早くお金を送って」——映っている人の顔も声も本物の息子そっくり。でも、偽物です。いま最も防ぎにくく、危険な詐欺手口の一つです。

\textbf{ステップ1：破綻を見抜く——AI 偽装によくある4つの弱点}\par

\begin{itemize}
  \item まばたきの頻度が不自然：人は1分間に15〜20回まばたきしますが、AI 顔すり替えではまばたきが少なすぎたり、リズムがぎこちなかったりしがちです。
  \item 口の動きと音声がずれる：口の形と聞こえる声の間に、0.数秒のズレがないか、よく見ましょう。
  \item 顔のふちがにじむ：顔の輪郭、耳のあたり、髪の生え際に「にじみ」やちらつきが出やすくなります。
  \item 光と影の向きが合わない：顔のハイライトや影が、背景の光の当たり方と食い違っていたり、横顔に向いたときに目鼻立ちがうまく「ついて回らなかったり」します。
\end{itemize}

\begin{tcbinfo}{ステップ2：3つの方法で相手が本物か確かめる}
\begin{itemize}
  \item 特定の動作をお願いする：相手に顔を向けさせたり、顔や口を手で覆わせたりして、口を開け閉めさせます。AI 顔すり替えは、顔が隠れたり大きく首を回したりするときに最も「ボロが出やすく」、顔がゆがんだり動きに追いつかなかったりします。
  \item 本人しか知らない質問をする：たとえば「この前、うちの犬をどこの病院に連れて行ったっけ？」といった嘘の質問を用意します——詐欺師なら答えられないか、口ごもるはずです。
  \item わかっている番号にかけ直す：一度ビデオ通話を切り、連絡先に登録してある番号へ直接かけ直します。詐欺師は本物の電話番号を操れないので、かけた瞬間に正体が分かります。
\end{itemize}

\end{tcbinfo}

\begin{tcbwarn}{ステップ3：騙されたと確かめたあとの通報の手順}
\begin{enumerate}
  \item まず送金を即停止：どんなに「急いで」と言われても一度ストップ。「早く振り込め、電話を切るな」と急かす相手は、ほぼ間違いなく詐欺師です。
  \item 証拠をすべて残す：ビデオ通話の画面を録画·スクリーンショットし、チャット履歴、相手のアカウント、送金の証憑を保存します——これらは技術鑑定の重要な材料です。
  \item 96110（詐欺被害相談専用ダイヤル）へ電話：「AI 顔すり替え／声まね」タイプの詐欺だと伝えると、オペレーターが緊急の支払い停止を始めてくれます。
  \item 110 に通報：証拠を一式そろえて派出所へ行き、事情聴取を受けます。AI による偽装技術が使われたことを特にはっきり伝えると、警察の立件と性質の判断に役立ちます。
  \item 知人·家族に知らせる：詐欺師があなたの偽造動画で他人を騙すおそれがあるので、モーメンツや家族グループであらかじめ注意を呼びかけましょう。
\end{enumerate}

\end{tcbwarn}

\booksection{シーン7：SIM スワップ攻撃（SIM 乗っ取り）——あなたの番号が「盗まれる」}

聞いたことがない人が多いのに、被害が非常に大きい攻撃です。SIM スワップ（SIM カード乗っ取り）とは、攻撃者があなたの個人情報を偽装してなりすまし、キャリアショップで新しい SIM カードを再発行することです。再発行に成功すると、あなたのスマホの電波は突然消え、攻撃者のスマホはその新しいカードで、あなた宛ての SMS 認証コードをすべて受け取れます——Alipay も銀行も WeChat も、すべての扉が大きく開いた状態になります。

\begin{tcbinfo}{どう気づく？SIM スワップの3つの予兆サイン}
\begin{itemize}
  \item スマホが突然「圏外」／「SIM カードがありません」と表示される。カードには触れていないのに、再起動しても直らない——多くの場合、これが乗っ取られた瞬間です。
  \item キャリアから「カード再発行が完了しました」「手続きが完了しました」という SMS が届く——しかし自分では何も申し込んでいない。
  \item 突然認証コードが届かなくなり、ほかの SMS も届かなくなる——SMS が攻撃者のカードへ横取りされているからです。
\end{itemize}

\end{tcbinfo}

\begin{tcbwarn}{防御策——SIM カードにも「鍵」をかける}
\begin{itemize}
  \item SIM カードの PIN を設定：［設定］＞［電話］＞［SIM PIN］から。有効にすると、スマホの再起動やカードの差し替えのたびに PIN の入力が必要になります。攻撃者がカードを手に入れても使えません。
  \item キャリアの二段階確認を有効にする：中国移動／聯通／電信のアプリにログインし、「セキュリティセンター」や「手続き時の確認」を探して、本人以外による再発行の SMS 通知と、手続き時の二重パスワード確認をオンにします。こうすれば、他人があなたのカードを再発行しようとしたとき、窓口の係員が追加の本人確認を求めてくれます。
  \item 重要なアカウントは SMS 認証だけにしない：銀行や Apple ID など重要度の高いアカウントは、SMS の認証コードではなく、認証アプリ（Google Authenticator、Microsoft Authenticator）やハードウェアキーを優先しましょう。
  \item 身分証のコピーには透かしを：手続きに必要なコピーには「○○の手続きにのみ使用」と書き添えます——攻撃者は、流出した身分証の情報であなたになりすますことが多いからです。
\end{itemize}

\end{tcbwarn}

\textbf{乗っ取られたあとの緊急対処——スピードがすべて}\par

\begin{enumerate}
  \item すぐにキャリアのサポートへ電話して利用停止：移動10086／聯通10010／電信10000。ただちに回線を止め、カード再発行の手続きを凍結するよう依頼し、攻撃者があなたの番号を使い続けられないようにします。
  \item 別のデバイスで重要アカウントにログイン：パソコンか家族のスマホから、銀行、Alipay、WeChat、Apple ID のパスワードをすぐに変更し、見知らぬデバイスを締め出します。
  \item 決済アカウントを凍結：Alipay は 95188 で利用停止、WeChat はセキュリティセンターで凍結、銀行はサポートにネットバンキングの凍結を依頼します。
  \item 身分証を持ってキャリアショップへ：本人が身分証を持参して直営店に行き、新しい SIM カードを再発行して、番号の利用権を取り戻します。
  \item 警察に通報：SIM スワップは「個人情報の侵害」にあたります。証拠（回線停止の時刻の記録、金銭的損失、再発行の記録）をそろえて被害を届け出ましょう。
  \item 総点検：番号が戻ったら、直近のログイン記録と入出金明細を一つずつ確認し、なりすましの形跡がないか確かめます。
\end{enumerate}

\begin{tcbwarn}{予防は治療に勝る——ふだんからの5つのこと}
\begin{enumerate}
  \item ロック画面のパスワード＋二段階認証——両方オンに
  \item 「デバイスを探す」が有効か確認
  \item SIM カードに PIN を設定（［設定］＞［電話］＞［SIM PIN］）——スマホをなくしても、他人にカードを使われません
  \item 出どころの分からないアプリを入れない、知らないリンクを開かない
  \item 大事なデータ（写真·連絡先·チャット履歴）を定期的にバックアップ
\end{enumerate}

\end{tcbwarn}

\begin{tcbpractice}{やってみよう：あなた専用の「緊急連絡カード」をつくる}
\begin{enumerate}
  \item 紙を1枚用意して、次の番号を書きましょう：キャリアのサポート（移動10086／聯通10010／電信10000）、銀行のサポート、96110 詐欺被害相談専用ダイヤル、110（警察通報）、家族の電話番号。
  \item Apple ID／各ブランドのアカウントの登録メールアドレスと、予備の連絡手段も書いておきます。
  \item この紙は財布に入れるか冷蔵庫に貼っておく——スマホの中にしまわないこと。何かあったときスマホが手元にないかもしれないからです。
  \item 今すぐチェック：「デバイスを探す」が有効か、SIM カードの PIN が設定済みか確かめましょう。
\end{enumerate}

\end{tcbpractice}

\begin{tcbsummary}{この章のまとめ}
\begin{itemize}
  \item スマホ紛失のゴールデン30分：借りたスマホで発信 遠隔ロック SIM 利用停止 決済凍結 知人·家族に連絡 警察へ通報
  \item 騙されたあとの72時間：送金を止める 証拠を残す 96110 に電話 通報 銀行に連絡 パスワード変更
  \item スマホの感染：セーフモード 怪しいアプリを削除 セキュリティスキャン 初期化（最終手段）
  \item アカウント乗っ取り：パスワード変更 デバイスを締め出す ログイン履歴を確認 二段階認証を有効化 連携情報を確認 決済設定を確認
  \item 脅迫メッセージの99.9\%は詐欺。返信も支払いもせず、スクリーンショットを撮ってブロック
  \item AI 顔すり替え／声まね詐欺：まばたき·口の動き·輪郭·光影の破綻をチェック。顔を向ける·口を覆う動作を頼み、本人しか知らない質問をする、登録済みの番号へかけ直す
  \item SIM スワップ乗っ取り：スマホが突然「圏外」になったら警戒。SIM PIN の設定、キャリアの二段階確認、重要アカウントは認証アプリへ切り替え
  \item ふだんの予防5か条：ロック画面＋2FA、デバイスを探す、SIM PIN、出どころ不明アプリを入れない、定期バックアップ
\end{itemize}

\end{tcbsummary}
